
\subsection{A-Levels}

A-levels are the terminal, standardized qualifications taken at the end of upper-secondary education (age 16–18) and form the principal basis for university admission (Figure~1). Applicants may apply to up to five courses via UCAS, typically by a January deadline; because the examinations are held after applications close, applications are supported by teacher-predicted grades. Universities then issue conditional offers (usually by May) that specify the grades required for admission. Applicants nominate a firm and an insurance choice; these choices establish the binding order in which offers are accepted.

Final placement is determined by the centrally awarded A-level grades. A-level scripts are marked by examination boards: markers are blind to candidates' identities and to application information, and letter grades are assigned according to standardised grade boundaries. Once actual results are released in August, UCAS assigns students according to the ranking they submitted and the conditions of the conditional offers (the firm choice is accepted if its conditions are met; otherwise the insurance choice may be used; failing both, the student can enter Clearing). Because grading and placement follow this centralised, rule-based procedure, teachers and schools have no formal discretion over the final marks or over the contractual placement outcome.

A-level scripts are marked centrally by examination boards; markers are blind to candidates' identities and demographic attributes. Letter grades are awarded according to grade boundaries that examiners set annually. Those boundaries are calibrated with reference to historical performance and other standardisation procedures so that the difficulty of achieving a given grade is broadly comparable across years. We confirm this property empirically by predicting A-level attainment from pre-treatment observables (here, GCSE mathematics and English) and comparing the conditional expectation of A-level outcomes across non-COVID years (\Cref{fig:event_study}): visual inspection of binned conditional expectation functions, pooled regressions with year × predicted-score interactions, and formal equivalence tests within narrow predicted-score bins reveal no meaningful shifts in the conditional means. Complementary checks — subject-by-exam-board splits — produce the same conclusion, supporting the claim that, conditional on prior attainment, A-level outcomes are stable across non-COVID cohorts.


Subject choice and course provision vary across the system. There are well over 100 distinct A-level subjects, and schools and sixth-form colleges decide which subjects to offer according to local demand and resources. Some pupils remain at their secondary school for sixth-form study, while others transfer to separate sixth-form or further-education colleges; this spatial and institutional dispersion affects which subjects pupils can take. Universities often specify required subjects for particular courses, and highly selective institutions tend to favour so-called ``facilitating'' subjects; independent schools frequently encourage uptake of these subjects and invest in resources (e.g. specialised equipment and teaching) to support them.

% Unique feature 1: Large importance in admission
The UK admission system puts a large weight on the outcomes of the qualifications that student undertakes for offering a placement to an applicant. Students submits their applications to up to 5 university courses before they sit their A-levels, as the exam dates of the actual exams are after after the deadline for university application. The deadlines for the application varies by university but most universities set their deadline near January. In place of the exam grades, students often ask their teachers or instructor at their attending school to provide a prediction of their grade that students can use to signal their academic competence. After receiving the applications with a predicted or past performance in a standardized exam, the universities reply to the offer with a grade requirement that the student must provide at the end of the application cycle. The deadline for universities to send their replies are usually near May.  After receiving the offers from the universities, students must choose which university they want to prioritize by choosing two universities that did not reject their application, and declining the offers from the other universities if they received any. Among the two universities that they choose, they must assign a rank, which dictates which university the students get assigned when the final grades are determined, due to the centralized admission system (UCAS) assigning students to their top rank school if they passed the grade requirement. After sitting the actual A-levels exams (in June) and receiving the results (around August), students are automatically assigned to a university based on the ranking they submitted beforehand. 

% Unique feature 2: subject choice
The subjects choices are broad, ranging from well studied subjects such as mathematics and English literature, to rare subjects such as Greek studies. On average, the number of unique subjects studies by student was near 100 during the sample period of this study. The subjects are usually connected to the subject that the students want to study in university, as universities set requirements on the subjects that students need to take in order to be assessed for admission. Competitive schools, such as the Russell Group schools, request students to take subjects from a set of subjects called \emph{facilitating subjects} so that they can assess the quality of the applicants more easily.  Since there are more than 100 A-levels in the UK, secondary schools choose which subjects are provided to their students. Independent schools often advise their pupils to take more facilitating subjects to increase their chance of getting accepted into a higher ranked school. In addition, many independent schools spend their capital on equipment to improve their pupils' performance in facilitating subjects, which attracts applications from parents. 

%Grading of the A-levels
The A-levels are graded by a centralized examiner that is blinded with any attributes of the students taking the exam. After the exams are graded, a letter grade is assigned to the score based on absolute performance of the student in the exam. However, the grade boundaries are set by a myriad of examinations, namely the examiners using historical data on the performance of previous year student and setting the boundaries of the grades so that the difficulties in achieving a letter grade is constant across multiple years. 

\subsection{University application system in the UK}
The UK employs a centralized clearing house to facilitate the matching between students and universities. Students can only apply up to 5 universities and they must state their preferences over their universities after receiving replies to their applications. The universities reply to applications before the date that students sit the A-levels. The reply content are broad from immediate acceptance and rejections, as well as conditional offers that state the grade requirement that students must pass for securing placements at the university program. Before sitting the A-levels, students must choose two university programs that they are willing to enroll, and discard the offers from the other three universities. The two chosen university programs are ranked which dictates the placement of the student after they receive the A-level results. The choices are called \emph{firm choice} and \emph{insurance choice} respectively. 


\subsection{COVID-19 pandemic}
The COVID-19 pandemic forced the UK government to cancel A levels and replace them with a teacher-based classification scheme. To prevent the spread of the virus, governments canceled in-person exams and replaced them with teacher-based assessment grading. Teachers of the applying pupils were asked to predict the scores of the students if the A-levels had taken place. Teachers resorted to their knowledge of the student and the students' past performance in their classroom. 

The cancellation of the A-levels were announced on March 18th 2020 for the 2020-2021 academic cohort and January 6th 2021 for the 2021-2022 cohort. In both years, students had already sent applications to their preferred university programs and received replies from universities that stated the outcomes. The outcomes included conditional offers, which states the grade requirements that students had to meet to secure a placement in the university program, unconditional acceptance into the program, and rejection.  

\subsection{Consequence of the method change}
The immediate consequence of the method change was a substantial amount of grade inflation (\Cref{fig:grade_histogram}). The number of students receiving A* or A increases by approximately 10 percentage points compared to pre-pandemic shares in 2020, and by 20 percentage points in 2021. Notably the extent of the grade inflation was substantially larger in the second year, as the 2021 cohort experienced a larger disruption in their preparation for the A-levels than the previous cohort. 

The inflated grades pushed students into universities, as universities could not control their inflow of students. Universities were legally obliged to accept students that passed their grade conditions stated in their conditional offers. \Cref{fig:CohortSizeGrowthAny} and \Cref{fig:CohortSizeGrowthFirm} shows how universities over accepted students in both 2020 and 2021. The expansion of the incoming class size was concentrated in schools with lower acceptance rate in the pre-pandemic years, as the cohort size nearly doubled in the most competitive program despite the programs having a larger class size than lower ranked programs. \Cref{fig:CohortSizeGrowthFirm} provides how the expansion in cohort size was driven by students choosing the university as their firm choice. Notably the number of university programs with acceptance rates lower than the 30th percentile rank increased their accepting size more than twice the pre-pandemic average. 

\subsection{The British Education Sector for Higher Education}
The UK has on average 3,300 university programs per year, matching with on average 570,000 students per a year. 3,500 secondary schools provide students with education to prepare for exams.

British secondary schools are mostly run publicly. 93\% of students participate in public schools and only 7\% in private schools ((\Cref{tab:balance})). Private schools, equally known as independent schools, run mainly on tuition collected from their pupils' parents. Past literature documents that attendants of independent schools are often children with rich parents. Henseke et al (2021) finds that 92 percent of the students attending independent schools belong to households in the top 10th decile group of the income distribution. 90\% of the students in independent schools proceed to higher education after graduation, with 4\% attending Oxford of Cambridge and 52\% attending one of the top 25th ranked universities. Students from independent schools apply to better universities more than students from public schools. Almost 85\% of the students apply to at least one university of the Russell group, while 45\% of the students from public schools apply. Independent school students apply less to lower-tier schools than students from public schools. Almost 40\% of the students apply to lower-tier schools, while 80\% of the students from public schools apply to at least one lower tier school. Independent school students receive more offer of admission from Russell Group schools than public schools. 80\% of students from independent schools receive at least one offer from a Russell Group school. 

There are multiple types of public school in the UK higher education sector (\Cref{tab:balance}). The most common type of schools are academy schools that consists of 35\% of the total number of schools. Academy schools provide education to pupils living within their local authority jurisdiction. As an exception, grammar schools select students based on their performance on a previous national exam (\emph{year 11}). 


%Table [] reports the simulation results for income sorting across universities by grading regime.  Column 1, 2, and 3 reports the correlation between university rank and the average parental income of students for incoming student cohort for standardized, non-standardized grading scheme for 2020 and 2021, respectively. The results indicate how the sorting pattern across university by income strengths for both 2020 and 2021 grading patterns. A one percent increase in the selectivity of a university, measured by the incoming cohorts average tariff scores, is associated with a 0.44 increase in the standard deviation of the parental income distribution.  The slope is statistically significantly larger than the magnitude of the coefficient with the standardized exam grading patterns. 

%Table [] reports the simulation results for the average academic competence of the incoming cohort measured by students GCSE math score. The non-standardized regime predicts that the average sorting pattern between students academic competence and the selectivity of the university strengths, with a one percent increase in the universities selectivity leading to a 1.04 standard deviation point reduction in the incoming cohorts average GCSE math score. The estimate for the standardized exam regime was -0.69, and it is statistically smaller in magnitude than the non-standardized regime. [Without endogenous application.] 

%Table [] reports the simulation results for the distribution of females among incoming across universities. Although the female students receive higher grades than male students in the non-standardized testing regime, the regression coefficient between female share and the selectivity of the university decreases. A one-percent increase in the selectivity of the university is associated with a 0.2 percentage point decrease in the share of female students with the non-standardized grading patterns. The estimate is -0.07 percentage point for the standardized grading regime.
